\documentclass{article}

\usepackage{tabularx}
\usepackage{booktabs}

\title{Problem Statement and Goals\\\progname}

\author{\authname}

\date{}

\input{../Comments.text}
\input{../Common.text}

\begin{document}

\maketitle

\begin{table}[hp]
\caption{Revision History} \label{TblRevisionHistory}
\begin{tabularx}{\textwidth}{llX}
\toprule
\textbf{Date} & \textbf{Developer(s)} & \textbf{Change}\\
\midrule
Date1 & Name(s) & Description of changes\\
Date2 & Name(s) & Description of changes\\
... & ... & ...\\
\bottomrule
\end{tabularx}
\end{table}

\section{Problem Statement}

\wss{You should check your problem statement with the
\href{https://smiths.github.io/capTemplate/Checklists/ProbState-Checklist.pdf}
{problem statement checklist}}. 

\wss{You can change the section headings, as long as you include the required
information.}

\subsection{Problem}

\wss{What problem are you trying to solve?  Why is it important? Do not talk
about how you will solve the problem, only what the problem is.  AI is rarely
part of the problem.  The problem is ``what'' you want to do, not ``how'' you
want to do it.  The problem statement should be written in a way that is
understandable to a non-technical audience.}

\subsection{Inputs and Outputs}

\wss{Characterize the problem in terms of ``high level'' inputs and outputs.  
Use abstraction so that you can avoid details.}

\subsection{Stakeholders}

\subsection{Environment}

\textbf{Hardware Environment:} The game requires a device that supports web browsers (e.g., computers, laptops, etc.). Additionally, standard input devices including keyboard and mouse are used to interact with the game.

\textbf{Software Environment:} The game is accessed through a web browser (e.g., Google Chrome, Microsoft Edge, etc.). An internet connection is required for access to the website as well as to update the online components of the game.

\section{Goals}

\section{Stretch Goals}

\wss{Goals you hope to get to, but not necessary.  They are also helpful if the
scope of the original project needs to be expanded.}

\section{Custom Documents (CDs)}

\wss{For CAS 741: State whether the project is a research project. This
designation, with the approval (or request) of the instructor, can be modified
over the course of the term.}

\wss{For SE Capstone: List your custom documents (CDs).  Potential CDs include
usability testing, code walkthroughs, user documentation, formal proof,
GenderMag personas, Design Thinking, etc.  (The full list is on the course
outline and in Lecture 02.) Normally the number of CDs will be two (one per
term).  Approval of the CDs will be part of the discussion with the instructor
for approving the project. The CDs, with the approval (or request) of the
instructor, can be modified over the course of the term.}

\newpage{}

\section*{Appendix --- Reflection}

\wss{Not required for CAS 741}

\input{../Reflection.text}

\begin{enumerate}
    \item What went well while writing this deliverable? 
    \item What pain points did you experience during this deliverable, and how
    did you resolve them?
    \item How did you and your team adjust the scope of your goals to ensure
    they are suitable for a Capstone project (not overly ambitious but also of
    appropriate complexity for a senior design project)?
\end{enumerate}  

\subsection*{Allen Tong}

\begin{enumerate}
    \item What went well for me was that the sections were split up clearly from the start, so I always knew what I was responsible for and didn't have to guess or wait on the rest of the team to figure out my part.
    \item One pain point for me was understanding the rules of the game well enough to explain clearly, since the rules document provided by Dr. Rapoport is technical and uses octal. We each went through the documents on our own, then held a group meeting with Dr. Rapoport to ask clarifying questions, which gave us a better understanding of how the game works. Another pain point was deciding which ideas should be goals and which should be stretch goals, since we had many possible features. We discussed this as a team and used the risk and complexity of each idea to decide.
    \item Dr. Smith told us that our game as originally described was too easy for an eight month senior capstone project for a team of five. To fix this, we added goals that take real engineering effort to implement beyond the base game, which include 3D visualization of the cube and an online competitive mode. This gave us enough complexity for a senior project without making the required goals depend on features we are unsure we can finish. 
\end{enumerate}

\subsection{Josh Engel}

\begin{enumerate}
    \item I feel as though the goal section was the easiest to create as we had already gone through multiple iterations of our goal to allow for the project to be accepted by Dr.smith.
    This meant less time had to be given to the development of our goals in the document but instead it was smooth just formalizing the goals and strech goals we had already created.
    \item The main point I experienced was simply communication. On my end I did not know any of my groupmates before capstone and so am not the most comfortable with reaching out on a consistent basis.
    With this in mind I was waiting for others to begin a conversation on issues we were experiencing and so went for periods of time with no knowledge. In the future I plan to reach out more consistently 
    whenever I am struggling with an issue to get help.
    \item We had to expand the scope of our project since the initial game concept was too simple for a capstone project. To do this we added the goals, 3D visualitzation, and competetive play, 
    as well as a strech goal of a family of games in different number bases.
\end{enumerate}  

\subsection{Nehan Munir}

\begin{enumerate}
    \item This deliverable was a lot easier than expected due to working with our supervisor, Dr. Rapoport. We had extensive meetrings with him, where he patiently explained how magic cubes work (which we weren't aware of before). He also had documents available to us that explains how these work if we forgot, alongside the rule set given to the users. He also gave us the eight preplanned levels, alongside an adequate scoring system. While this would restrict the design space for us, it was useful to wrap our head around the project idea.
    \item We did struggle with figuring out who the problem statement was for. Since we were making a game, we didn't know if the problem statement should be written fo the stakeholders (which are likely to be highschoolers) or for what we want to accomplish during the capstone (which will be more about the innovations that this project will be incorporating). Luckily, a discussion with the TA cleared up our misunderstanding.
    \item As described and written by Dr. Rapoport, the levels that he wanted to make can clearly be done in a less than a month. While we did discuss adding a family of games, another way to expand the scope is to add 3D visualization. Looking at how Dr. Rapoport and other examples found online, they were always using a 2D model, which isn't the best representation of the 8x8x8 cube, especially if we want to the user to be able to see triagonals. 
\end{enumerate}  

\subsection{Amos}

\begin{enumerate}
    \item I would say team cooperation went well for this deliverable. Starting the project with a new team often means that I am unsure if our communication styles and times would match upt but thankfully they did and it allowed us to be consistent with meetings and discussions so far.
    \item My main pain point was understanding the actual project itself as even after our first meeting with our supervisor I was very lost on what the game was exactly, we had a few meetings amongst ourselves and I started heavily documentign what was said so I could review it and finally in our second meeting I was able to adequately represent my confusions to the supervisor and along with the documents he provided beforehand get and understandding of magic 8.
    \item We added a few new features some we consider stretch goals and walked through them with our supervisor and TA to confirm they would add to the project meaningfully and be sufficiently complex.
\end{enumerate} 

\end{document}